% !TeX root = ../main.tex
% Chapter 6 -- Conclusion (Conclusao Geral)
% Draft 1 (2026-07-23). Written from NORTH_STAR section 6 Ch.6 beats (synced with all
% signed-off additions, AVAL rounds 1-3) + the D1 capacity-baseline record
% (storyline/audit/capacity_baseline_experiment.md section 5).
% Obeys WRITING_LAW.md + GLOSSARY.md. Claim/number ledger: 6_citations.md (same folder).
% [AUTHOR DECISION pending in section 6.2: prominence of the capacity-matched baseline
%  paragraph -- it is included per the D1 licensing contract (suppression is not an
%  option once run), but its length and placement are the author's call.]


\chapter{Conclusion}
\label{ch:conclusion}

This dissertation asked whether multitask learning helps predict the next category and
the next region of a visit, and which design choices determine the answer. Three studies
addressed this question in sequence. The first developed a joint model that did not
consistently outperform its dedicated counterparts (Chapter~\ref{ch:cbic}). The second
kept the model architecture fixed, changed the input representation, and identified that
representation as the main bottleneck in the tested configuration
(Chapter~\ref{ch:courb}). Those two studies paired next-category prediction with a static
category classification task, and the pair changed in the final study.
% AUT-32 (PENDENCIAS §4, ruling "Opção B"). The author noticed that this opening does not mention the
% static category classification task, and the omission is real. Option B was chosen over widening the
% research question, and the distinction matters: the question stated here is the SAME question the
% NORTH_STAR §1 and §1.2 state, about the next category and the next region. Adding the static task to
% the QUESTION would make this chapter declare a research question that neither of those two approved
% places declares, so the change would propagate rather than stay a wording fix. The clause above
% instead records the static task as HISTORY of the first two studies, which is what it is.
% Consistent with the same treatment elsewhere: the static task is present in the Ch.3/Ch.4 recaps and
% in limitation 6, and absent from the research question and from the future work. See AUT-35 (c) and
% AUT-36, which turn on the same task-pair change.
% Probe: R13-aut32.
The third study combined a check-in-level representation with a redesigned sharing
topology (Chapter~\ref{ch:mobiwac}). Under this design, one joint model replaces two
dedicated single-task models on both tasks. On the region task it outperforms
them at Texas and California, the two datasets with the largest region vocabularies, and at
the other four it stays within a two-point margin registered before any result was read
(statistical non-inferiority, assessed with two one-sided tests). On the category task it
outperforms them at Florida, and the five remaining differences are equivalent to zero within
half a point.
% [NEEDS SIGN-OFF 30 | CLOSED] Author confirmed the scope as written: category at all six, region
% at four (Istanbul, Florida, California, Texas), TOST non-inferiority within two points at
% Alabama and Arizona, Arizona not upgraded.


\section{Contributions by chapter}
\label{sec:conc:contributions}

Chapter~\ref{ch:cbic} presents MTLnet, the first joint model developed in this research.
It combined static category classification with next-category prediction and shared its
lower layers across the two tasks. Both tasks read a place-level graph embedding. In
this configuration, the joint model cost more to train and did not consistently
outperform the dedicated single-task models on either task. The chapter contributes
this negative finding and an analysis of three possible causes: task dissimilarity and
the resulting risk of negative transfer, an input representation that may be
insufficient for both tasks, and an overly rigid sharing topology. The finding is
limited to a place-level input under hard sharing. The next two studies examine what
happens when these conditions change.

Chapter~\ref{ch:courb} kept the MTLnet architecture fixed but replaced its single
64-dimensional place embedding with separate spatial, temporal, and categorical
encoders. On the static task, category macro-F1 rose by 20.2 to 22.0 percentage points
across the three states tested. This increase exceeded those obtained from the
architectural variations evaluated in the first study.
% [round6, D-02] The 20.2-22.0 point gain is the STATIC task's (4_courb/results.tex:92), whose
% input determines its target by construction (Appendix B), so it does not by itself diagnose the
% arc; the sequential task, clean in both published chapters, is what Ch.5 goes on to test. The
% number stays as the published chapter's own audited figure, now carrying its task and
% qualification instead of standing unlabelled as the whole arc's diagnosis.
% [NEEDS SIGN-OFF 31 | CLOSED] Author approved the two added frame-prose sentences; no number
% changed.
% [NEEDS SIGN-OFF 32 | CLOSED] Author: no longer applicable given the current state of the text.

This gain requires two qualifications: First, the static task classifies a place from
that place's own representation, so the input already determines the target. The gain
therefore provides no evidence about the sequential task. Second, the comparison is not
width-matched: the decomposed input has 192 dimensions, whereas the place embedding has
64. Chapter~\ref{ch:courb} therefore calls for an equal-dimension control to separate the
semantic contribution of the encoders from the effect of the additional width.

The sequential task provides the relevant diagnosis because its input does not
determine its target. The decomposed encoders are more effective in most of the tested
cases. Because the sharing topology remained fixed, this result identifies the input
representation as the main bottleneck in that configuration.

Chapter~\ref{ch:mobiwac} responds to this diagnosis with Check2HGI, a check-in-level
representation that assigns a distinct vector to each visit rather than reusing one
vector for every visit to the same place. The joint model exchanges information between
two task-specific streams through cross-attention, while a private spatial path serves
the region task. The evaluation used user-disjoint cross-validation, with twenty fitted
models per configuration: four seeds, each drawing its own five-fold user partition.
Paired tests used the four per-seed means.
% CORRECTED 2026-08-04, the same finding the author caught in Appendix A. This said "four seeds over one
% fixed set of five folds", which is false: src/data/folds.py passes the run's seed into
% StratifiedGroupKFold(random_state=...), so the partition moves with the seed. Evidence and the full code
% path: science/fold_partition_and_seeds.md. This was the LAST live occurrence in the .tex sources; the
% sweep that found it is recorded in §9 of that document.
% "Paired tests used the four per-seed means" is UNCHANGED and is what keeps this sentence coherent: the
% tests never required fold k of one seed to be fold k of another. Probe: R13-foldseed9.

Under this protocol, the joint model outperforms the dedicated single-task models on the
region task at Texas ($+1.21$ Acc@10) and California ($+1.06$), and at Istanbul, Alabama,
Arizona, and Florida it stays within the registered two-point margin. All four are small
deficits, the largest being Alabama at $-0.87$; a test in the reverse direction, corrected across
the six region comparisons, resolves three of them. They are reported as differences inside the
margin rather than as ties. On
the category task it outperforms the dedicated model at Florida ($+0.19$ macro-F1); the
other five are equivalent to zero within half a point. The equivalence margin was registered for
the region axis only, so the category bound is read off the intervals rather than chosen in advance.
Non-inferiority was the registered primary analysis for the region task; the
two superiority findings are secondary results outside that analysis plan. Further analyses
clarify what these results do and do not establish. They isolate the contribution of the
input representation, test parameter count as an alternative explanation for the next-category
gains, and identify a possible relation between problem scale and the next-region gains.

Beginning with input representation, Section~\ref{sec:mobiwac:results-part1} provides the first
explanation through a controlled comparison between Check2HGI and HGI. With the dedicated
model, training configuration, windowing, and folds held fixed, Check2HGI produces higher
next-category macro-F1 than the place embedding HGI across the six evaluated datasets.
This result establishes the input representation as one condition that affects
performance in this setting. It also supports testing Check2HGI in other mobility
prediction architectures, although its benefit in those architectures has not yet been
evaluated.

Model capacity offers a competing explanation for the next-category results, which this
dissertation attributes to the check-in-level representation.
Appendix~\ref{apx:parameter-count-control} tests it directly: the dedicated category model is
widened to the joint model's parameter budget, holding the representation, the folds, and the
training protocol fixed. At Alabama, where the control ran the full four-seed
protocol, the wider model scores $0.53$ macro-F1 points below the narrower one it replaces, and the
direction is unanimous across all twenty folds. Capacity is therefore not the variable next-category
prediction responds to at that dataset, which leaves the representation as the explanation the
evidence supports. A single-seed reading at California is consistent with it.

For next-region prediction, the results point to problem scale as another possible
condition. The two datasets where the joint model outperforms the dedicated model are
Texas and California, the two states with the largest number of regions, while the four
datasets with smaller region vocabularies sit inside the margin. This grouping suggests
that joint learning becomes more useful when the model must choose among more regions.
Two observations keep it a suggestion. The ordering does not hold inside the pair, since
California has more regions than Texas and a slightly smaller gain, and states with more
regions also tend to have more check-ins. A controlled
experiment is therefore required to separate the effect of region count from the amount
of data and other differences between the states. Problem scale remains a possible
condition, not an established cause.

% [2026-08-06, AUTHOR REWRITE] The findings on representation, capacity, and problem scale remain
% separate in the preceding section. The consolidated answer follows three movements: (1) the
% question and final result; (2) the sequence from the Chapter 3 null result, through the Chapter 4
% diagnosis, to the Chapter 5 design; and (3) the incomplete initial premise and conditional lesson.
% The first result remains valid for its configuration. The correction applies to a broad reading of
% that result, not to the result itself. The MTL expectation is attributed internally to
% Section~\ref{sec:fund:mtl}, where it is already sourced. The Check2HGI implication is limited to a
% candidate for testing with other architectures; transfer to them is explicitly unmeasured. The
% distinction between diagnosis and resolution is carried by the paragraph sequence.


\section{The consolidated answer}
\label{sec:conc:answer}

The research question asks whether multitask learning helps next-category and
next-region prediction, and what the answer depends on. For the first part, the evidence
shows that it helps under the final configuration. The sequence that led to this
configuration also identifies the conditions behind that answer. Identifying these
conditions is the main finding of this dissertation.

The explanation begins with Chapter~\ref{ch:cbic}, which used a place-level embedding and
naive hard sharing. Its joint model did not consistently outperform the two dedicated
models, leaving three possible explanations: the tasks might be too different, the input
might be insufficient, or the sharing topology might be too restrictive.
Chapter~\ref{ch:courb} separated one of these explanations by
keeping the architecture fixed and changing the input. Performance then improved,
including on the sequential task in most states. Because the architecture did not
change, this comparison identifies the input representation as the main bottleneck in
that configuration. This diagnosis set the direction for the final study.

Chapter~\ref{ch:mobiwac} responded to this diagnosis in two steps. First, it introduced
Check2HGI, which replaced one vector per place with one vector per visit. Second, the
joint architecture replaced naive hard sharing with task-specific streams that exchange
information through cross-attention, while the region output retained a private spatial
path.

Together, these changes define the final configuration evaluated in this dissertation.
They also explain why the answer is conditional. The initial premise, explained in
Section~\ref{sec:fund:mtl}, was that training related tasks together would improve
generalization through a shared representation. Across the three studies, task
relatedness and joint training were not sufficient by themselves.
The controlled comparisons identify the input representation as one condition,
while the results suggest that the architecture and dataset scale may matter, although their
individual effects were not fully isolated. The answer therefore depends on the design and
evaluation conditions, not on joint training alone.

These findings support a conditional rather than universal conclusion.
The consolidated answer is therefore not that multitask learning always helps.
It is that multitask learning helps next-category and next-region prediction under
the final design and evaluation protocol developed in this dissertation.

\section{Limitations}
\label{sec:conc:limitations}

Six limitations bound the scope of these conclusions.

\begin{enumerate}
    % [round8, 2026-07-30, PENDENCIAS_RESOLVIDOS 5.6 (arquivado 2026-07-30)] BOTH DATES, because they are different
    % facts
    %     and one of them
    % was about to be printed alone. The author's decision was "busque pelo que o artigo original cita e
    % vamos usar isso em ambos", and a track did exactly that: opened cho2011gowalla first-hand (authors'
    % posted PDF, cs.stanford.edu, Section 2 p.2) and put Feb 2009 - Oct 2010 here. Correct as a statement
    % ABOUT THAT PAPER, and wrong as a statement about this work's data, which is what a limitation headed
    % "Data vintage" is read as.
    % MEASURED on the five parquet files this work consumes (columns: local_time; command below):
    %   Alabama    2009-03-18 .. 2011-07-27   n=113,846
    %   Arizona    2009-03-26 .. 2011-07-04   n=236,450
    %   Florida    2009-03-13 .. 2011-08-11   n=1,407,034
    %   Texas      2009-01-21 .. 2011-08-16   n=4,089,892
    %   California 2009-01-24 .. 2011-08-14   n=3,171,380
    %   union 2009-01-21 .. 2011-08-16 -- TEN MONTHS past the window the paper states.
    %   python3 -c "import pandas as pd,glob,os
    %   for f in ['Alabama','Arizona','Florida','Texas','California']:
    %       t=pd.to_datetime(pd.read_parquet(f'/Users/vitor/Desktop/mestrado/data/checkins_by_state/{f}.parquet')
    %       ['local_time'],errors='coerce')
    %       print(f, t.min(), t.max(), len(t))"
    % The author's premise in 5.6 ("ambos usaram o mesmo recorte, nao houve diferenca") is therefore not
    % what the files show. Not a defect in his decision -- he was ruling on which SOURCE to cite -- but the
    % sentence carried both dates for two days; it now carries only the measured one.
    % [round9c, PENDENCIAS_RESOLVIDOS 5.6b (arquivado 2026-08-02), 2026-07-30] SIGN-OFF MARKER REMOVED, question
    % answered. It asked which
    % window to print and offered the cut: "If he prefers only the paper's window, the clause after the
    % comma is the one to cut." He chose the OTHER direction, only the measured window, so the clause that
    % came out is the paper's. Nothing is left for him to decide here, and his instruction is explicit:
    % "O need sing-off assim como os demais ja resolvidos que estao no latex pode ser removidos nao
    % precisam fica lá." The reasoning is preserved in _round9/45_author_rulings.md, per the second half of
    % the same instruction ("Se quider documentar isso tem que ser em algum lugar do src_util").
    \item \textbf{Data vintage.} The five state datasets come from
    Gowalla~\cite{cho2011gowalla}, and the extraction used here spans January 2009
    to August 2011 across the five states. The Istanbul
    dataset~\cite{wongso2025massivesteps} is not appreciably more recent for most of its
    volume: its check-ins fall in two separate periods, 2012 to 2013 and 2017 to 2018,
    with none in between, and roughly seven in ten belong to the earlier period.
    Mobility patterns, place inventories,
    and check-in behavior have changed since.
    % AUT-35 (a) (PENDENCIAS §4, ruling "Vamos aplica o A,B e commentar a limitação 6, com as
    % alterações do C", plus the author's follow-up choice of the two-block form on 2026-08-04).
    % "August 2011" is LOAD-BEARING: probe R8-vintage requires that literal string, so the Gowalla
    % window must not be reworded when this item is edited.
    % THE AUTHOR'S PREMISE FOR THIS EDIT WAS REFUTED, and the sentence above is what survived.
    % He believed the Istanbul check-ins were 2017-2018 data and that Massive-STEPS is the most modern
    % public benchmark. Neither holds as stated:
    %   - the check-ins span 2012-04-03 to 2018-10-19 in TWO blocks with nothing in 2014, 2015 or 2016.
    %     Of the 462,615 Istanbul check-ins this work models, 327,242 (70.7%) fall in 2012-2013 and
    %     135,373 (29.3%) in 2017-2018. MEASURED, not quoted: the benchmark paper gives no per-city
    %     window anywhere. Command: python3 src_utils/_round13/aut35a_window.py datetime
    %     output/check2hgi/istanbul/chrono_split/split_assignment.parquet, corroborated by
    %     data/massive_steps_istanbul/parse_report.json (field datetime_range_local). The instrument
    %     was blindness-checked first against the Gowalla files, where it returns 2009/2010/2011 and
    %     nothing else, which is the window this same item states.
    %   - "most modern public dataset" fails on Yelp, whose open dataset carries check-ins to
    %     2022-01-19 (measured by streaming checkin.json). The narrow form that does survive concerns
    %     only openly redistributable, city-partitioned check-in benchmarks used in this literature,
    %     and it is a claim about the Massive-STEPS COLLECTION rather than the Istanbul split. That
    %     claim is not made in this document.
    % The prose says "roughly seven in ten" rather than 70.7%, because a per-dataset percentage of a
    % third party's data does not belong in a limitation item; the exact figure lives here with its
    % command. Full evidence and the source ledger: src_utils/_round13/70_massivesteps_validation.md.
    % Probes: R13-aut35a, and R8-vintage still guards the Gowalla string.
    % [round9c, PENDENCIAS_RESOLVIDOS 5.6b (arquivado 2026-08-02), author decision, 2026-07-30] ONE WINDOW, THE
    % MEASURED ONE. His words:
    % "Vamos usar so as datas que encontramos no database de 2009 a 2011, pode omitir que no artigo eles
    % comentam que e de 2009 a 20010." So the clause reporting what the dataset paper states (February
    % 2009 to October 2010) is REMOVED, and the citation stays on the dataset itself, which is what it is
    % for. "since either date" became "since", since there is now one date range.
    % The measurement that produced the printed span is unchanged and is preserved in full in the
    % provenance block above: five parquet files, union January 21 2009 to August 16 2011, with the
    % command that reads them. The paper's own window is not lost either -- it is recorded in
    % src_utils/_round9/45_author_rulings.md under this item, which is where he asked documentation of
    % this kind to live ("Se quider documentar isso tem que ser em algum lugar do src_util").
    % [COD/PENDENCIAS_RESOLVIDOS 5.6 (arquivado 2026-07-30), round8, 2026-07-30] "collected between 2009 and 2011" ->
    %     "collected between
    % February 2009 and October 2010", with the citation added. Author decision, PENDENCIAS_RESOLVIDOS 5.6 (arquivado
    %     2026-07-30): "Busque
    % pelo que o artigo original cita e vamos usar isso em ambos. Inclusive ambos usaram o mesmo recorte
    % nao houve diferenca."
    % SOURCED THIS SESSION, first hand. Cho, Myers and Leskovec, "Friendship and Mobility: User Movement
    % in Location-Based Social Networks", KDD 2011, DOI 10.1145/2020408.2020579. The publisher copy is
    % closed (Unpaywall, Semantic Scholar, PMC, Elsevier and Springer all returned no open copy); the
    % authors' own posted PDF was opened at cs.stanford.edu/people/jure/pubs/mobile-kdd11.pdf, 9 pages.
    % Section 2, "CHARACTERISTICS OF CHECK-INS", page 2, states: "We collected all the public check-in
    % data between Feb. 2009 and Oct. 2010 for Gowalla and Apr. 2008 to Oct. 2010 for Brightkite. The
    % total number of check-ins for Gowalla is 6.4 million". Only the window is used here; the 6.4 million
    % figure is NOT quoted in the text, because it counts the whole SNAP collection and not the five-state
    % subset this work uses.
    % CHAPTER 4 NEEDS NO EDIT, which is why this is a one-file change rather than two. Its published prose
    % at 4_courb/conclusion.tex:20 already reads "collected between February 2009 and October 2010", the
    % paper's window verbatim, so the two chapters now agree and no published text is touched. Checked
    % against the Portuguese record and the English translation, ../CoUrb_2026/src/sections/conclusion.tex
    % and src_en/sections/conclusion.tex: both carry the same window.
    % A MEASURED FACT THIS DOES NOT ERASE, and it is the reason for the marker below. The provenance
    % comment at 5_mobiwac/05_setup.tex:17-21 records that the data this work actually consumes is the
    % figshare CC0 category-annotated dump (36,001,959 check-ins), whose date range was MEASURED on the
    % parquet on 2026-07-09 as 2009-01-21 to 2011-08-16, and it states that the SNAP/cho2011 dump is not
    % the data source. That measurement is where "2009 and 2011" came from, so it was not a careless
    % number, and it is left in place unaltered. The sentence above is therefore written as a statement
    % about the GOWALLA COLLECTION as its own paper reports it, not as a claim about the span of the five
    % parquet files, which is the only form in which the author's instruction and the measurement can both
    % stand. PENDENCIAS_RESOLVIDOS 5.6 (arquivado 2026-07-30)'s own warning applies and is respected: do not "correct" a
    %     measured number to
    % match an inherited one.
    % [NEEDS SIGN-OFF 36 | CLOSED] Author confirmed both studies used the same extract.
    \item \textbf{Taxonomy coarseness.} Next-category prediction uses seven top-level
    classes. A finer taxonomy may change the effect of joint training.
    \item \textbf{Transductive representation.} Check2HGI is trained on the check-in
    graph of each dataset, so it cannot represent unseen places or users without
    retraining. The place-level representation it builds
    on~\cite{huang2023hgi} is trained the same way, and methods that learn aggregators
    instead of per-node vectors~\cite{hamilton2017graphsage} are the alternative this
    constraint points toward.
    % AUT-35 (b) (PENDENCIAS §4, same ruling as (a) and (c)). The author wanted to add that the
    % transductive constraint "affects several approaches in the literature". That is a claim about
    % OTHER PEOPLE'S work, and AGENT_GUARDRAILS §1 forbids it as a bare generality: it needs a located
    % citation, and the citation must come BEFORE the claim rather than be appended to it, which is the
    % order used here. Both poles were already in references.bib with resolvable identifiers:
    % huang2023hgi (10.1016/j.isprsjprs.2022.11.021), whose HGI is the place embedding this work's
    % baseline uses and is trained per dataset, and hamilton2017graphsage, whose title is literally
    % "Inductive Representation Learning on Large Graphs".
    % WHAT IS DELIBERATELY NOT WRITTEN: no count, no "several", no "many methods". The sentence names
    % two identified works and stops, because "several approaches" is a survey claim and no survey was
    % read in first person to support it. A broader statement needs a survey sentence opened and quoted,
    % not an adverb.
    % NOTE FOR WHOEVER EDITS THIS ITEM NEXT: this edit was applied twice. The first application was
    % lost when a parallel editing lane rewrote this file at 02:49 while the round was in progress, and
    % the loss was caught only because probe R13-aut35b failed. The probe is the reason it did not ship
    % missing.
    % Probe: R13-aut35b.
    \item \textbf{No next-place task.} The experiments do not predict the exact next
    POI, and their conclusions apply only to next category and next region.
    \item \textbf{Geographic coverage.} Outside the United States, the evidence rests on
    a single city, Istanbul.
    \item \textbf{The task-pair confound.} The task pair changed together with the
    representation and the sharing topology. Chapters~\ref{ch:cbic}
    and~\ref{ch:courb} paired static category classification with next-category
    prediction, while Chapter~\ref{ch:mobiwac} pairs two sequential targets,
    next-category and next-region prediction. No single controlled ablation
    separates the representation-and-topology change from the task-pair change in
    the final result. Chapter~\ref{ch:courb} is the fixed-pair control for the
    diagnosis. The additional parameter-count experiment in
    Appendix~\ref{apx:parameter-count-control} tests only whether a larger dedicated
    category model can close the gap. It covers Alabama under the four-seed
    protocol and California at one seed, one width point per dataset, width scaling
    rather than depth, and a smaller hyperparameter search than the original
    dedicated model. The greater similarity of
    the final pair may still contribute to the size of the improvement. The
    ablation that would separate the two changes,
    running static category classification under the check-in-level representation, is
    not clean under that representation: the category of the visited place is an input
    feature of a check-in node, so that target would be partly readable from its own
    input. This follows from the design of the representation and was not measured, so
    the confound is bounded by the fixed-pair control rather than removed.
    % AUT-35 (c) (PENDENCIAS §4, ruling "commentar a limitação 6, com as alterações do C").
    % THIS ITEM WAS COMMENTED OUT IN THE WORKING TREE AND WAS RESTORED ON THE AUTHOR'S INSTRUCTION,
    % 2026-08-04. It is live at HEAD 9813ac21 and all twelve lines had been prefixed with % by a
    % parallel editing lane, with no provenance comment, which is not this repository's convention. Two
    % consequences made it a defect rather than a choice: the section says "Six limitations" and would
    % have rendered five, and §6.4 ties future-work items to limitation NUMBERS, so shortening the
    % enumeration leaves limitation~6 pointing at nothing. CORRECTED after a reviewer finding: this
    % comment previously said the removal "silently repoints limitation~1 through limitation~4", which is
    % false and I stated it as fact after hedging it correctly when I asked the author. This item is the
    % LAST in the enumerate, so removing it renumbers nothing above it; the tags 1 through 5 keep their
    % referents and only the tag for this item dangles. The load-bearing reason to restore it is the one
    % that does hold: the section opens with "Six limitations" and would have rendered five.
    % The item also carries a
    % registered sign-off from 2026-07-22, which AGENT_GUARDRAILS requires be renewed before removal.
    % THE ADDED SENTENCES ARE MARKED AS INFERENCE, and that marking is load-bearing. The leakage
    % reasoning is ANALYTIC, derived from the representation's construction (the category one-hot is an
    % input feature of the check-in node, preprocess.py:615-642), and NO experiment measured it. The
    % text therefore says "follows from the design" and "was not measured" rather than reporting a
    % finding. Do not upgrade those clauses to a result.
    % WHY THIS STRENGTHENS THE LIMITATION RATHER THAN WEAKENING IT, which was the author's own reading:
    % if the ablation that would resolve the confound cannot be run cleanly, then the confound is not
    % cleanly resolvable, which is a stronger bound on the conclusions than "no ablation was run".
    % What the argument DOES invalidate is the future-work item that proposed running that ablation.
    % See AUT-36 for the coupled change: NORTH_STAR §6 requires each future-work item to be tied 1:1 to
    % a limitation, so this limitation and that item move together.
    % Probes: R13-aut35c, R13-lim6live.
    % [GATE FIX C-1] "the size of the win" -> "the size of the improvement" (idiom law).
    % [NEEDS SIGN-OFF 37 | CLOSED] Author: already qualified in the live prose above (category
    % task, two of six datasets, one width point per dataset, width scaling rather than depth).
\end{enumerate}


\section{Future work}
\label{sec:conc:future}

Future work follows directly from these limitations. Newer and denser traces would
test the conclusions beyond the Gowalla vintage (limitation~1), and finer-grained
taxonomies would test them beyond the seven-class division (limitation~2). Because the
check-in-level representation is trained on a fixed graph, an inductive variant of it,
one that embeds unseen places and users without retraining, would remove the
transductive constraint and support deployment in growing cities. Two further changes to
the same representation would test how much of its behavior follows from how it is
assembled rather than from the hierarchy: its place-vector table is initialized from a
pretrained table and then held near it during training, so an ablation that varies that
coupling would separate the contribution of the initialization, and a hypergraph
formulation, in which one edge joins the several visits of a session rather than only
consecutive pairs, would test whether the pairwise graph is the right structure for a
visit history (limitation~3).
% AUT-36 (PENDENCIAS §4, ruling "Seguimos com até quatro frases novas ... Sobre o proximo lugar podemos
% adicionar ele no trabalho futuro relacionado ao uso do check2hgi para outros contextos. Esse eu faço
% questão de estar de alguma forma. Quanto ao resto adicione os que você analisou e cabem estar lá").
% THREE items were added, and they were attached to EXISTING limitations rather than creating new ones.
% (His ruling allowed up to four. A first draft of this edit did add a fourth, the cascade-tuning and
% soft-sharing item under limitation 6, and that version was replaced: limitation 6 already hosts the
% task-pair item, so tagging a second item to it broke the 1:1 rule this comment goes on to explain.
% The fourth item is therefore recorded as dropped below, not as added.)
% That choice is the whole structural point: NORTH_STAR §6 requires future work tied 1:1 to
% limitations, and §6.3 opens with "Six limitations", so inventing limitations to host new items would
% have changed the count sentence and silently repointed every limitation~N tag in this section.
%   - the initialization coupling and the hypergraph formulation attach to limitation 3, which is about
%     how the representation is built, and they are the two of his seven items that were ABSENT from
%     the whole tree (a hypergraph search returned zero hits across all .tex files).
%   - the next-place item carries his mechanism, and it is the one he insisted on ("faço questão").
%     TENSE DISCIPLINE IS LOAD-BEARING HERE: "would require a change to the input pipeline and one
%     additional output" is conditional. His own phrasing, "do jeito que está hoje já conseguimos
%     usar", must NEVER be carried over: the present tense would claim a capability, and next place is
%     a GATED scope exclusion (GLOSSARY §1.1; Definition 2.9 in 2_fundamentals.tex states that no
%     chapter reports a result for it, guarded by probes R12-fplace and R12-fplace2). His framing of
%     the item as Check2HGI reaching other contexts is what the clause carries.
%   - the cascade-tuning and soft-sharing item was DROPPED rather than squeezed in, because both
%     candidate hosts (limitation 6 and a new limitation) were unavailable: limitation 6 already hosts
%     the task-pair item, and adding a limitation changes the count. Recorded rather than silently
%     omitted, so it can be revisited if he authorizes a seventh limitation.
% Two of his seven items were ALREADY PRESENT and needed nothing: finer taxonomies (limitation 2) and
% further non-United-States cities (limitation 5).
% Probes: R13-aut36a, R13-aut36b.

Adding the exact next place as a third target would extend the joint model beyond the
two properties predicted in this work. The cascade architectures reviewed in
Chapter~\ref{ch:mobiwac} suggest that category and region predictions can provide
additional structure for this task, and reusing the existing check-in-level
representation would require a change to the input pipeline and one additional output
rather than a new representation, which is also the route by which this representation
would serve contexts other than the two studied here (limitation~4).
Further cities outside the United States would widen the geographic base
(limitation~5).
Isolating the effect of changing the task pair needs a static target
that the check-in-level representation does not already carry as an input feature, since
running the earlier pair directly under this representation is not a clean comparison
(limitation~6).

% AUT-36 (PENDENCIAS §4) COUPLED CHANGE, and it is not optional. The item here previously proposed
% "training the Chapter 5 joint model on the Chapter 3 task pair, under the check-in-level
% representation", which is precisely the ablation that AUT-35 (c) establishes is not clean: the
% category of the visited place is an input feature of a check-in node, so the static target would be
% partly readable from its own input. Leaving both texts unchanged would have the limitation explain
% why an experiment cannot be run cleanly while the future work proposed running it.
% NORTH_STAR §6 requires every future-work item to be tied 1:1 to a limitation, so the tie to
% limitation 6 is preserved and only the proposal changes. It now names the condition a valid version
% of the comparison would have to meet instead of naming an experiment that cannot meet it.
% Probe: R13-aut36lim6.


\section{Final remarks}
\label{sec:conc:final}

The practical output of this dissertation is a single model that predicts two
properties of the next visit in one forward pass: what kind of place the user will
visit and in which part of the city. The methodological contribution is the sequence
of evidence that led to this model: a published negative result, the diagnosis that
identified the input representation as the main bottleneck, and a solution designed
from that diagnosis. The negative result was not an obstacle to the contribution. It
was its first half.
